\documentclass[11pt]{article}

\usepackage[margin=1in]{geometry}
\usepackage{amsmath}
\usepackage{amssymb}
\usepackage{booktabs}
\usepackage{enumitem}
\usepackage{hyperref}

\title{Batched Monotone Tracing}
\author{}
\date{}

\begin{document}

\maketitle

\section{Goal}

The batched monotone tracing method estimates the fidelity-$0.5$ contour in the
two-dimensional randomized benchmarking configuration space
\[
    (n, r) =
    (\texttt{n\_gates}, \texttt{ratio\_2\_qb\_gates}),
\]
where $n$ is the circuit gate count and $r$ is the two-qubit gate ratio.
For each fixed ratio $r$, the method searches for a local bracket
\[
    p(n_{\mathrm{lo}}, r) > 0.5,
    \qquad
    p(n_{\mathrm{hi}}, r) < 0.5,
\]
and records the midpoint as the crossing estimate
\[
    n_*(r) \approx
    2 \operatorname{round}\left(
        \frac{n_{\mathrm{lo}} + n_{\mathrm{hi}}}{4}
    \right).
\]

Unlike a pure bisection method, this implementation probes several candidate
gate counts in one stitched batch. This is intended to amortize the fixed
submission cost while keeping each fixed-ratio search local to the currently
predicted boundary.

\section{Code Organization}

The implementation is split across:

\begin{itemize}[leftmargin=*]
    \item \texttt{batched\_monotone\_tracing.py}: settings, backend setup,
    ratio scheduling, run orchestration, saving, plotting, and scoring.
    \item \texttt{batched\_tracing\_helpers/hints.py}: line, surface, and
    analytic contour predictions; conversion of predictions into search
    windows.
    \item \texttt{batched\_tracing\_helpers/windows.py}: batched probing,
    posterior confidence tests, bracket extraction, per-ratio budget handling,
    and the fixed-ratio search routine.
    \item \texttt{batched\_tracing\_helpers/constraints.py}: monotonicity
    constraints and local jump guards.
\end{itemize}

\section{High-Level Algorithm}

Let $B$ be the remaining global HQC budget and let $C$ be the list of accepted
crossings. The main loop is:

\begin{enumerate}[leftmargin=*]
    \item Initialize the backend, random number generator, RMB object, data
    store, and global budget using \texttt{start\_run}.
    \item Choose one or more ratio sweeps from \texttt{trace\_direction}.
    The default direction is downward: start at the upper ratio bound and trace
    toward the lower ratio bound.
    \item For each ratio $r$ in the sweep, call \texttt{search\_one\_ratio}.
    \item If a confident local bracket is found, convert it to a midpoint
    crossing and append it to $C$.
    \item If no bracket is found, print a failure message and move to the next
    ratio.
    \item Stop when the ratio sweep ends or the global budget is exhausted.
    \item Save crossings, print the experiment summary, optionally print the
    score, and optionally plot the results.
\end{enumerate}

\section{Ratio Scheduling}

\subsection{Start Ratio}

The start ratio is set by \texttt{start\_ratio}. If this is \texttt{None}, the
default is determined from \texttt{trace\_direction}:
\[
\begin{cases}
    r_{\max}, & \texttt{trace\_direction = "down"},\\
    r_{\min}, & \texttt{trace\_direction = "up"},\\
    (r_{\min} + r_{\max}) / 2, & \texttt{trace\_direction = "both"}.
\end{cases}
\]

\subsection{Fixed and Adaptive Steps}

If \texttt{adaptive\_ratio\_step} is false, ratios are separated by
\texttt{ratio\_step}. If it is true, the step is adjusted from the remaining
span and the remaining HQC budget:
\[
    \Delta r =
    \operatorname{clip}
    \left(
        \frac{|r_{\mathrm{stop}} - r|}
             {\max(1, \lfloor B / \texttt{target\_hqc\_per\_ratio}\rfloor)},
        \texttt{min\_ratio\_step},
        \texttt{max\_ratio\_step}
    \right).
\]
This makes the sweep spend fewer ratio slices when the budget is already low.

\section{Predictions and Search Windows}

At each ratio, the method chooses a vertical gate-count window $[n_{\min}, n_{\max}]$
to probe. The prediction logic uses three possible sources.

\subsection{Analytic Prediction}

The analytic Lindblad contour is evaluated through
\texttt{analytic\_gate\_counts}. This returns a predicted crossing
\[
    \hat n_{\mathrm{analytic}}(r).
\]
It is used as a cheap prior, especially before enough measurements exist to
fit a surface.

\subsection{Line-Fit Prediction}

Once at least two accepted crossings exist, the method fits a line in inverse
gate count:
\[
    \frac{1}{n_*(r)} \approx a r + b.
\]
The predicted gate count is then
\[
    \hat n_{\mathrm{line}}(r) = \frac{1}{a r + b},
\]
provided the denominator is positive.

\subsection{Surface Prediction}

If \texttt{use\_surface\_bracket\_hint} is true and at least
\texttt{surface\_min\_configs} measured configurations exist, the method fits
the shared monotone fidelity surface using
\texttt{try\_fit\_monotone\_fidelity\_surface}. It evaluates the fitted
surface along the vertical slice at fixed ratio $r$ and finds the first sign
change of
\[
    \hat p(n, r) - 0.5.
\]
Linear interpolation between adjacent grid points gives
$\hat n_{\mathrm{surface}}(r)$.

\section{Initial Anchor Search}

The first accepted crossing is treated specially because no previous contour
points are available.

\begin{enumerate}[leftmargin=*]
    \item Try to predict the initial crossing using the surface hint.
    This usually only works if the run starts with existing data.
    \item If the surface hint is unavailable, use the analytic prediction.
    \item If neither prediction is available, fall back to a broad low-to-high
    window
    \[
        [n_{\mathrm{low}},\;
        \texttt{initial\_bracket\_fraction} \cdot n_{\mathrm{max}}].
    \]
    \item If a prediction $\hat n$ exists, center a window on it:
    \[
        h =
        \max\left(
            2,\;
            \texttt{initial\_prediction\_window\_fraction}\cdot \hat n,\;
            \frac{\texttt{initial\_prediction\_min\_width}}{2}
        \right),
    \]
    and probe $[\hat n - h, \hat n + h]$, clamped to
    \texttt{n\_gates\_bounds}.
    \item Submit \texttt{initial\_batch\_points} candidate gate counts, each
    with \texttt{initial\_batch\_shots} shots.
\end{enumerate}

For broad windows, candidate gate counts are spaced geometrically using
\texttt{low\_to\_high\_growth\_factor}. For prediction-centered windows, they
are linearly spaced.

\section{Tracing After an Anchor}

Once one or more crossings have been accepted, the method uses the previous
crossings to form a wider safe bracket and then narrows it to a focused batch
window.

\subsection{Wide Bracket}

The function \texttt{next\_bracket} first chooses predictions according to
\texttt{bracket\_hint}, which may be \texttt{"line"}, \texttt{"surface"}, or
\texttt{"both"}. If predictions are available, their median is used as the
center. The half-width is
\[
    w =
    \max\left(
        \texttt{bracket\_half\_width\_fraction}\cdot \hat n,\;
        \max_i |\hat n - \hat n_i|,\;
        2
    \right).
\]
The bracket is then enlarged using the most recent crossing. When tracing
downward in ratio, the high side may expand to
\[
    n_{\mathrm{last}} \cdot \texttt{trace\_gate\_growth},
\]
and the low side may include
\[
    n_{\mathrm{last}} \cdot \texttt{trace\_gate\_shrink}.
\]
The analogous inverse scaling is used when tracing upward.

If no prediction is available, the method falls back to a bracket based only on
the last accepted crossing and \texttt{trace\_gate\_growth} or
\texttt{trace\_gate\_shrink}.

\subsection{Monotonic Bounds}

The wide bracket is intersected with bounds implied by nearby accepted
crossings. The assumed contour is monotone: moving to lower two-qubit ratio
should require the same or more total gates to reach fidelity $0.5$.

Only the nearest accepted crossing above and below the current ratio are used.
This avoids overconstraining the search with distant noisy anchors. A slack
band is added:
\[
    s =
    \max(
        \texttt{monotonic\_min\_slack\_gates},\;
        \texttt{monotonic\_slack\_fraction}\cdot n_{\mathrm{anchor}}
    ).
\]

When tracing downward, an additional growth floor can be imposed from the
analytic contour:
\[
    g =
    \max\left(
        \texttt{downward\_growth\_floor\_min\_gates},\;
        \texttt{downward\_growth\_floor\_fraction}
        \left[
            \hat n_{\mathrm{analytic}}(r)
            -
            \hat n_{\mathrm{analytic}}(r_{\mathrm{anchor}})
        \right]
    \right).
\]
This prevents the trace from staying vertically flat for too long when the
analytic model expects gate count to increase.

After at least \texttt{trace\_gate\_jump\_min\_anchors} accepted crossings,
local jump guards also limit unusually large jumps. The allowed jump from an
anchor is scaled by
\[
    \min\left(
        \texttt{trace\_gate\_jump\_step\_scale\_cap},\;
        \max\left(1,\frac{|r-r_{\mathrm{anchor}}|}{\texttt{ratio\_step}}\right)
    \right).
\]
The growth side is controlled by \texttt{max\_trace\_gate\_growth\_fraction};
the shrink side is controlled by \texttt{max\_trace\_gate\_shrink\_fraction}.

\subsection{Focused Batch Window}

The function \texttt{trace\_batch\_window} collects all available predictions:
surface, line, and analytic. It chooses a prediction quantile rather than
always taking the median:
\[
    q =
    \begin{cases}
        \texttt{downward\_prediction\_quantile}, & r < r_{\mathrm{last}},\\
        \texttt{upward\_prediction\_quantile}, & r > r_{\mathrm{last}},\\
        0.5, & \text{otherwise}.
    \end{cases}
\]
This predicted center is converted into a narrower window:
\[
    h =
    \max\left(
        2,\;
        \texttt{trace\_batch\_window\_fraction}\cdot \hat n,\;
        \frac{\texttt{trace\_batch\_min\_width}}{2}
    \right).
\]
The final probed window is $[\hat n-h,\hat n+h]$ clamped inside both the wide
bracket and the monotonic bounds.

\section{Batched Probing}

For a chosen ratio and gate-count window, the method creates
\texttt{MeasurementRequest} objects for a small set of candidate gate counts.

\begin{itemize}[leftmargin=*]
    \item Initial-anchor windows use \texttt{initial\_batch\_points} and
    \texttt{initial\_batch\_shots}.
    \item Trace windows use \texttt{trace\_batch\_points} and
    \texttt{trace\_batch\_shots}.
    \item Broad windows use geometric spacing.
    \item Prediction-centered windows use linear spacing.
\end{itemize}

Requests are packed into stitched batches subject to
\texttt{max\_cost\_per\_run}. Before a batch is submitted, its HQC cost is
checked against both the global budget and the current ratio budget.

\section{Budgeting}

The first anchor search may use the full remaining global budget. After at
least one crossing has been accepted, each ratio is capped by
\texttt{max\_hqc\_per\_ratio}, if that parameter is not \texttt{None}. The
effective spendable amount for a batch is
\[
    \min(B_{\mathrm{global}}, B_{\mathrm{ratio}}).
\]

This keeps later fixed-ratio searches from consuming the entire remaining
budget when a local window fails to bracket cleanly.

\section{Bracket Acceptance}

After the batch is measured, configurations are scanned in increasing gate
count. For each measured configuration, the code computes the posterior mean
and posterior probability that the true fidelity is above $0.5$:
\[
    P(p(n,r) > 0.5 \mid \mathrm{data}).
\]

A point is treated as confidently above the boundary when
\[
    P(p(n,r) > 0.5 \mid \mathrm{data})
    \geq
    \texttt{decision\_confidence}.
\]
A point is treated as confidently below the boundary when
\[
    1 -
    P(p(n,r) > 0.5 \mid \mathrm{data})
    \geq
    \texttt{decision\_confidence}.
\]

The first adjacent transition from a confidently above point to a confidently
below point is accepted as the local bracket. If no such transition appears,
the fixed-ratio search returns no crossing and the main loop moves on.

\section{Crossing Validation}

Given an accepted bracket $[n_{\mathrm{lo}}, n_{\mathrm{hi}}]$, the recorded
crossing is the even midpoint. Before appending it, the method checks the
monotonicity constraints again. If the midpoint falls outside the allowed
monotonic bounds, it is rejected and not added to the trace.

\section{Key Parameters}

\begin{table}[h]
\centering
\begin{tabular}{lll}
\toprule
Parameter & Default & Role \\
\midrule
\texttt{ratio\_step} & 0.1 & Base spacing between ratio slices. \\
\texttt{trace\_direction} & \texttt{"down"} & Direction of contour tracing. \\
\texttt{adaptive\_ratio\_step} & \texttt{True} & Enables budget-aware ratio spacing. \\
\texttt{min\_ratio\_step} & 0.1 & Smallest adaptive ratio step. \\
\texttt{max\_ratio\_step} & 0.2 & Largest adaptive ratio step. \\
\texttt{target\_hqc\_per\_ratio} & 35.0 & Expected spend used by adaptive stepping. \\
\texttt{decision\_confidence} & 0.70 & Posterior confidence threshold for brackets. \\
\texttt{max\_hqc\_per\_ratio} & 35.0 & Per-ratio cap after first accepted crossing. \\
\texttt{initial\_batch\_points} & 5 & Number of initial-anchor gate probes. \\
\texttt{initial\_batch\_shots} & 3 & Shots per initial-anchor probe. \\
\texttt{trace\_batch\_points} & 5 & Number of tracing gate probes. \\
\texttt{trace\_batch\_shots} & 3 & Shots per tracing probe. \\
\texttt{bracket\_hint} & \texttt{"line"} & Prediction sources for wide brackets. \\
\texttt{use\_surface\_bracket\_hint} & \texttt{True} & Enables fitted-surface predictions. \\
\texttt{surface\_min\_configs} & 12 & Minimum measured configs before surface fitting. \\
\texttt{initial\_prediction\_window\_fraction} & 0.35 & Initial window width as fraction of prediction. \\
\texttt{initial\_prediction\_min\_width} & 80 & Minimum initial prediction window width. \\
\texttt{trace\_batch\_window\_fraction} & 0.14 & Trace window width as fraction of prediction. \\
\texttt{trace\_batch\_min\_width} & 40 & Minimum trace prediction window width. \\
\texttt{downward\_prediction\_quantile} & 0.75 & Prediction quantile when tracing downward. \\
\texttt{upward\_prediction\_quantile} & 0.25 & Prediction quantile when tracing upward. \\
\texttt{trace\_gate\_growth} & 2.4 & Wide-bracket growth allowance. \\
\texttt{trace\_gate\_shrink} & 0.65 & Wide-bracket shrink allowance. \\
\texttt{enforce\_ratio\_monotonicity} & \texttt{True} & Enables monotonic contour guards. \\
\texttt{monotonic\_slack\_fraction} & 0.65 & Slack around neighbouring anchors. \\
\texttt{max\_trace\_gate\_growth\_fraction} & 0.9 & Local growth jump guard. \\
\texttt{max\_trace\_gate\_shrink\_fraction} & 0.7 & Local shrink jump guard. \\
\texttt{downward\_growth\_floor\_fraction} & 0.45 & Fraction of analytic downward growth enforced. \\
\texttt{downward\_growth\_floor\_min\_gates} & 24 & Minimum downward growth floor. \\
\bottomrule
\end{tabular}
\end{table}

\section{Practical Interpretation}

The method is designed to balance three competing pressures:

\begin{itemize}[leftmargin=*]
    \item Batching improves HQC efficiency, so the method probes several nearby
    gate counts at once rather than repeatedly submitting one circuit.
    \item Prediction-centered windows avoid spending most of the budget far
    from the boundary.
    \item Monotonic bounds and jump guards reduce catastrophic excursions, but
    they must retain enough slack to avoid locking onto a bad early anchor.
\end{itemize}

The most sensitive parameters are usually the window widths, the posterior
confidence threshold, the per-ratio budget, and the monotonic growth guards.
Too narrow a window can miss the contour; too wide a window wastes the benefit
of batching. Too weak a monotonic guard allows vertical drops or large jumps;
too strong a guard can reject legitimate contour movement after a noisy anchor.

\end{document}
